\section{Two public clocks for counting computations}\label{sec:gap}

A batch and its description can both be exponential in a power of $n$.
The enumeration bound of Theorem~\ref{thm:closed} measures work in
terms of the whole network. The diagonal argument needs a different
bound. A public computation must specify the bits shared between
processors, while each processor can spend polynomial time in $n$
checking its own bits. We will distribute a long public computation
among private jobs whose inputs fit on individual processors.
We adapt the distribution of random access traces in
\cite[Section~\ext{G}{app:trace}]{oaG}.

The public machine is a universal random access machine (RAM) on
$w$-bit words. A tick acts on a bounded number of words, can address
memory indirectly, and can request one binary guess. Its fixed finite
instruction set includes Boolean operations, comparison, shifts and
addition. Multiplication, division, parsing and indexing take
$\operatorname{poly}(w)$ ticks on word-sized data. The machine starts
with a description of at most $R$ words, including its program, input
and any advice. Pad its reserved $R$ words with zeros and initialize
all other cells to zero. Instructions can equivalently
be fetched from read-only memory. We count paths that halt with output
one within the prescribed cap.

For small positive rational constants $b,\rho$, let
\[
    n^\rho\le \log R\le 3n^\rho,\qquad
    w=\lceil 20\log(nR)\rceil,\qquad
    T_-=\lfloor R n^{1-b}\rfloor,\quad
    T_+=\lceil R n^{1+b}\rceil .
\]
Our arrays have size $R n^{O(1)}$ with small absolute exponents, or
length $o(R)$. Counters of $O(w)$ bits use a bounded number of words.
Given $(n,R)$ explicitly, the programs compute the clocks and other
rounded powers with fixed rational exponents by integer comparisons or
bit lengths in $\operatorname{poly}(w)$ time.

At the smaller public cap $T_-$, the RAM may issue \emph{private jobs}.
Each job supplies at most $\lfloor n^{1/2}\rfloor$ bits to a fixed
multitape machine that can make binary guesses and has $n^{20}$ steps
per path. Its only input is the job string. The public cost includes one
tick per job and one per bit issued. A \emph{completed path} consists
of a public path and one private path for each job it issues. It
succeeds when the public path and all those private paths accept.
To use a quantity computed in private, the public machine guesses
its value and includes the claim in the job for verification.

\begin{theorem}[Clock gap]\label{thm:gap}
Fix the compressor identities of Section~\ref{sec:compr}. Choose
$0<b<1/20$ sufficiently small for these identities, then $0<\rho<b$
sufficiently small. The following statements hold for all sufficiently
large $n$, with the private machine fixed in (1).
\begin{enumerate}
\item \emph{Realizing the smaller clock.} For any fixed private machine,
$R$ copies of $U_n^*$ realize the number of completed successes at cap
$T_-$. Pin some slots to bits, pair the rest, and sum over the common
bit in each pair. The equality holds over $\mathbb Z$. From the size
parameters and the public description, a uniform RAM generates all
pairs and literal pins in $Rn\operatorname{poly}(w)$ ticks, with
polynomial exponent independent of $b,\rho$. The same pairing works
for every description of at most $R$ words at the given $(n,R)$,
with its literal pins set from the padded description.

\item \emph{Simulating the larger clock.} For a RAM computation at cap
$T_+$ without private jobs, a uniform procedure with a fixed private
machine uses $o(T_-)$ public ticks with jobs. Each completed success
has a monomial weight in the compressor scalars. Its public path
records the exponents, in $O(w)$ bits each. Over any commutative
unital ring satisfying the compressor identities, the sum of these
weights equals the simulated acceptance count in the ring. The
procedure takes any simulated description of at most $R$ words,
including one already in public memory. Its fixed compressor templates
use only the names of the scalars.
\end{enumerate}
\end{theorem}

The first construction uses only Boolean predicates and does not
require port admission. In the second, compression leaves enough
public time for additional $\operatorname{poly}(w)$ work per original
processor. Both diagonal arguments will call the first construction on
a supplied description that includes the actual program and advice.
For this reason its generator must reproduce the literal pins as well
as the pairings.

\subsection{Encoding a random access trace}

Both constructions use the same encoding of a RAM trace as a trivalent
Boolean network.
For a tick cap $U$ between $T_-$ and $T_+\operatorname{poly}(w)$,
it has one satisfying assignment per accepting path. Specified final
memory reads may be tested against supplied bits or exposed as outputs.
We use the guess-and-check representation of nondeterministic
computations \cite{Cook} and a sorting-based memory check
\cite{TinyRAM}.

Record initialization writes of the description. Split each tick into
a bounded number of microsteps making at most one memory access apiece,
including instruction fetches. Guess read values and any requested
nondeterministic bit, check the register transition, and record every
probe. Sort the initialization, probe and final-read records by address
and time, breaking any ties in a fixed order. Within each address,
adjacent records propagate the current value starting at zero. A write
replaces that value and a read must return it. Pin input literals and
acceptance, force unused guesses to zero, and pad a halted computation
uniquely. All auxiliary values, including the outputs of comparisons
and the current-value fields, are forced on a given path. Replacing
gates by their Boolean relations and distributing repeated bits by
binary or ternary equality gives the asserted multiplicity.

Use $O(U)$ rows, rounded up to a power of two, allowing $O(U)$ final
reads. Put one initialization, probe or output record in each occupied
row, with successive microsteps in consecutive rows. The contacts of a
row with other rows have a useful regularity. In each layer of
Batcher's bitonic sorting network \cite{Batcher}, a comparator's partner
row is obtained by XOR with a fixed mask. Consecutive rows have XOR
masks consisting of a suffix of ones. All transitions, memory checks
and comparisons therefore use only $\operatorname{poly}(w)$ types
of links per row. The comparisons, transition circuits and equality
gates within each row also have size $\operatorname{poly}(w)$.

Index a gate and each of its slots by row, phase (including comparator
layer), local operation and bit number. The bitonic recursion gives
the partner rows and comparator directions. Thus a row's size
parameters, index and incident literal bits suffice to generate all
its gates, pins, links and designated outputs in $\operatorname{poly}(w)$
time. The whole network costs $O(U\operatorname{poly}(w))$ to generate.
Its links are determined before any memory addresses are guessed.
All word-cost exponents here are independent of the program, $b$ and
$\rho$.

\subsection{Wiring the smaller clock}
\begin{proof}[Proof of Theorem~\ref{thm:gap}(1)]
The wiring has two tasks. It checks a public trace and sends each
private job in that trace to one processor. To supply the jobs in
fixed positions, augment the public computation with a stream in a
separate address space inaccessible to the simulated program. Write the stream in consecutive chunks of
$J=\lfloor n/8\rfloor$ bits. Each job record has a start bit, a
fixed-width length and the data bits. Before starting a job, open a
new chunk if the largest allowed record would not fit. Zero-pad the
unused part of each chunk. For large $n$ a maximum record uses less
than half a chunk, so each closed chunk, apart from the last one, is
at least half full. No job crosses a chunk boundary.

Reserve space for the worst case, allowing $O(w)$ record expansion
per charged job bit or job, a factor of two for padding, and rounding
to a whole chunk. The begin, append-bit and end instructions implement
the stream and reject malformed jobs. Checking the public cost cap
and record lengths, writing the stream, and reading every reserved bit
in address order on halt take $T_-\operatorname{poly}(w)$ ticks.
This control is fixed for all public descriptions. Use the trivalent
trace network at this enlarged cap, with the final stream reads as
outputs.

Wire each chunk's $J$ outputs to one $U_n^*$ whose remaining slots
specify a fixed program, $n$ if needed, and padding. The program
executes its jobs in order, accepting exactly when all accept. A
successful path selects one accepting private path for each job.
A chunk contains at most $n$ jobs of at most $n^{20}$ private steps
each. Even scanning the used tape at every simulated private step
keeps the joint computation within the processor's $n^{100}$ steps.

Pack the trivalent vertices into further processors. A fixed checking
program reads, for each assigned vertex, its truth table and at most
three incident bits. A padded list with its length specified costs a
constant number of bits per vertex apart from its header. Add trivalent
vertices as needed to connect literal pins and chunk bits. Wire the
two slots of each link, including when both belong to checks on a
single processor. Its program thus needs only the
individual checks.

The chunks and the packed vertices together use
\[
           O\bigl(1+T_-\operatorname{poly}(w)/n\bigr)
\]
processors. For fixed $b$, take $\rho$ small enough that this is less
than $R$ and pin the remaining processors to return one. Generating
the trace, packing it in order and listing its pins and pairs costs
$Rn\operatorname{poly}(w)$ ticks. On contraction, a valid trace has
one extension to the public checking bits. The chunk processors count
all independent choices of accepting private paths. Their product
gives precisely the number of completions of that public path.
\end{proof}

\subsection{Compressing the larger clock}
\begin{proof}[Proof of Theorem~\ref{thm:gap}(2)]
Set $U=T_+$. We put its verification trace on the cyclic butterfly
of Theorem~\ref{thm:closed}, using the fixed XOR links to generate the
wiring within the smaller public budget \cite[Section~\ext{G}{app:trace}]{oaG}.

\emph{Routing the rows.}
Fix $a=1/20$, and group consecutive rows into buckets of
$S=2^{\lfloor a\log n\rfloor}$ rows. The low row bits give the offset
in the bucket and the high bits give its lane. Put one bucket at
cursor zero on each lane. The two arcs to the next cursor position
either stay on the lane or flip its active coordinate. Add empty
parallel slots to reach the compressor's fixed degree. To route a
link with a prescribed inter-bucket XOR mask, traverse the cursor
cycle, flipping each coordinate
specified by the mask when its turn comes. At every cursor step
this induces a bijection of lanes. The coordinates already flipped
recover the source lane from a relay's lane.

Reserve slots on each arc by mask, exchange type, low source row
bits and bit number. A relay equates successive bits on an occupied
route and pins surplus slots to zero. Thus every satisfying trace
has a unique extension to the routed network. Its number of positions
$P$ and maximum arc width $D$ satisfy
\[
              P=O((U/S)w^{O(1)}),\qquad D=O(Sw^{O(1)}).
\]
The slot ranges give all capacities, arc identifiers, lane displacements
and bucket ranges in $P\operatorname{poly}(w)$ public work. They
require no scan of the $S$ low row indices of each bucket.
Description words enter only as literals of initialization rows. On
given arguments a bucket or relay factor counts locally guessed
assignments that pass the row predicates, wire equalities and padding
checks. Generating row descriptions and making these checks costs
$O(S\operatorname{poly}(w))$ per local random access path. At a relay
the source lane and type schedule identify occupied slots independently
of the description words. The same mask, type and row identifiers
identify the slots at their destinations.

\emph{Choosing the shared bits.}
Apply $t=\lceil z\log n\rceil$ blocks from Section~\ref{sec:compr}.
Choose $z>0$ rational and small enough that $C^t\le S^{1/4}$ for all
large $n$. The available $\log(U/S)+O(1)$ lane coordinates suffice.
The public machine generates the initial descriptors and reads the
guesses for a history in their canonical encoding. It checks admission
and constructs the recursive descriptors. Position and capacity
indices have $O(w)$ bits, so this costs $P\operatorname{poly}(w)$ ticks
per path. Guess the surviving bits once per arc, a total of at most
\[
              2^{-t}O(P(D+1))\le U n^{-z}\operatorname{poly}(w).
\]

\emph{Issuing the jobs.}
For each final position, issue its final arguments and the recursive
layout of its local contraction. Follow the recorded preimages to
generate slices, ties and internal field identifiers, using local
indices for each job. For example, write a new local index into an
array at the original identifier upon visiting it. Each occurrence
takes $\operatorname{poly}(w)$ public work. The layout gives each
initial bucket or relay by its identifier and size parameters. Supply
also the literal description words in the intersection of its bucket
range with the initialization segment. The public machine visits
actual literals and skips padding. The hidden field bits will be
guessed in private.

For each compressor scalar, guess a fixed-length exponent for each
position and supply it to that job as a claim. Sum these exponents in
public counters. Each path uses only polynomially many map entries
and polynomial monomial degrees in $C^t,D,w$ per position, and there
are at most $P$ positions. Canonical encodings of $O(w)$ bits
therefore suffice for every local exponent and its global total.
Each job has length at most
\[
                 (C^t)^2(D+S+1)\operatorname{poly}(w)<n^{1/4}
\]
once $\rho$ is small enough. Allowing for all steps, including issuing
each job, the public cost is
\begin{equation}\label{eq:dump}
                 (C^t P+R+U n^{-z})\operatorname{poly}(w).
\end{equation}
The first term pays for $O(C^t)$ layout records per final position
and the exponent claims. The second pays for the description words,
since each initial bucket moves intact to exactly one final position.
The third pays for the final arc bits supplied at their two endpoints.
The polynomial overhead covers addressing, guesses and admission checks,
including rejected paths.

\emph{Checking the private counts.}
Each private machine guesses the hidden fields, checks local map
patterns, and runs the path tests for its moved initial factors. It
accepts only if the map entries are allowed and their monomial
exponents sum to the supplied claims. There are at most $C^t$
occurrences and $(C^t)^2 O(D+1)$ field bits in one job. Slices and
ties can be checked bit by bit. For a map, let $m$ be the largest field length in its
motif. The selector scans in
\cite[proof of Theorem~\ext{G}{thm:conversion}]{oaG} test the entry
and determine its monomial exponents in $(1+m)^2\operatorname{poly}(w)$
steps.

For a fixed $q$, the bound $B=(C^t)^2(D+S+1)w^q$ covers both the
combined length of a job's expanded records and fields and the
largest field length in any of its motifs. A private RAM can store
the local records in arrays indexed within the job. Charge each occurrence
$B^2$ steps for scans and $O(S)$ for any original row group it holds,
with fixed polynomial factors in $w$. Since $S\le n^a$ and
$C^t\le S^{1/4}$, the resulting bound is
$C^t(B^2+S)w^{O(1)}\le n^{13a/4}w^{O(1)}<n^2$ for small $\rho$. Scanning lists
on tapes incurs at most quadratic simulation overhead, leaving room
for universal interpretation within $n^{20}$ private steps. A fixed
private program implements these tests. The public machine discards
nonadmitted histories before issuing jobs, and the checks can reject
malformed data within the caps.

Fix a history and all final arc bits. Write $\boldsymbol\alpha$ for
the list of compressor scalars and $N_j(e)$ for the number of accepting
private paths at position $j$ when its exponent claim is $e$. The
local contraction at $j$ is $\sum_e N_j(e)\boldsymbol\alpha^e$.
For publicly claimed exponents $(e_j)_j$ there are
$\prod_j N_j(e_j)$ completed successes, each with public weight
$\boldsymbol\alpha^{\sum_j e_j}$. Summing these contributions gives
the product of the local contractions. Summing over arc bits and
admitted histories now gives the original acceptance count by the
compressor identity.

It remains to fit the public cost into the smaller clock.
Since $C^t\le S^{1/4}$ and $U=T_+$, Equation~\eqref{eq:dump} is
\[
             O\bigl(U(n^{-3a/4}+n^{-z}+n^{-1-b})w^{O(1)}\bigr).
\]
Choose $b>0$ sufficiently small after $a,z$, for example
$2b<\min\{a/2,z\}/2$. The required saving is
$T_-/U=\Theta(n^{-2b})$. Since $w=\Theta(n^\rho)$, decreasing $\rho$
makes the displayed cost $o(T_-)$ and meets the wiring, job-length
and private-time bounds above. The finitely many exponents on $w$
in these estimates are fixed before $b,\rho$ are chosen.
\end{proof}
